\documentclass{situationist-preprint}

\title{Kuhn on Translation}
\msid{SCI-06-b}
\msfamily{SCI-06}
\msdiscipline{Philosophy of Science}
\mstype{Companion manuscript}
\mskeywords{Kuhn; incommensurability; translation; interpretation; theory choice}

\begin{document}
\maketitle

\begin{abstract}
In his later work Kuhn restated incommensurability as untranslatability: no translation preserves the content of a small set of interdefined terms across a revolution. We argue that Kuhn's own account, taken at its word, supports the conclusion of SCI-06, that successive paradigms are partially translatable and that theory choice remains rational across revolutions. Kuhn's local incommensurability leaves most of the vocabulary translatable; his distinction between translation and interpretation concedes that the untranslatable terms can be learned and their claims compared; and the comparison he describes is made in shared vocabulary by shared values. Untranslatability of a local cluster is compatible with the rational comparison of the theories that contain it.
\end{abstract}

\section{Kuhn's restatement}\label{s:kuhn}

In \emph{The Structure of Scientific Revolutions} incommensurability was a relation of many kinds, between problems, standards and ways of seeing \citep{Kuhn1970}. In later papers Kuhn narrowed it to a linguistic thesis \citep{Kuhn1983}. Two theories are incommensurable if there is no language, neutral or otherwise, into which both can be translated without loss. The incommensurability is \emph{local}: most terms common to the two theories function the same way in both, and only a small group of interdefined terms, such as \enquote{mass} and \enquote{force} in Newtonian mechanics, or \enquote{element}, \enquote{principle} and \enquote{phlogiston} in eighteenth-century chemistry, resists translation. The terms of such a group are learned together, and their meanings cannot be fixed one at a time.

Kuhn also distinguished translation from interpretation. A translator substitutes expressions of one language for those of another; an interpreter, like a historian reading an old text, learns the other language and can then understand its claims, even where no substitution is available. Incommensurability, on Kuhn's later view, blocks translation of the local cluster, not interpretation.

\section{Partial translatability}\label{s:partial}

\begin{proposition}\label{p:partial}
If incommensurability is local, the theories on either side of a revolution are partially translatable: all their sentences not containing terms of the cluster are translatable homophonically, and sentences containing them are translatable token by token where a token's reference can be fixed.
\end{proposition}

\begin{proof}
The first part is Kuhn's own characterization of local incommensurability. The second follows from the argument of \citet{Kitcher1978}: a term of the cluster, such as \enquote{dephlogisticated air}, has a heterogeneous reference potential, referring on some occasions through a description that fails and on others through the substance a chemist has in hand. When Priestley described the properties of dephlogisticated air produced in a particular experiment, the token referred to oxygen, and the sentence is translatable by a sentence about oxygen. What fails is a uniform translation of the type, not the translation of each token.
\end{proof}

Proposition~\ref{p:partial} is all that the first part of the conclusion of SCI-06 requires. It does not claim that every sentence of the earlier theory has an equivalent in the later one, only that the untranslatable remainder is a bounded part of the vocabulary, surrounded by translatable claims.

\begin{proposition}\label{p:choice}
Local untranslatability does not prevent the rational comparison of the theories that contain it.
\end{proposition}

\begin{proof}
Comparison requires that the theories' observational claims and the values by which they are judged be shared. The observational claims are expressed largely in the translatable vocabulary of Proposition~\ref{p:partial}. The values are, on Kuhn's own account, shared across revolutions: accuracy, consistency, scope, simplicity and fruitfulness, applied differently by different scientists but common to the community on both sides \citep{Kuhn1977}. And the untranslatable claims are accessible to interpretation, so that a chemist who has learned both languages can assess what each theory says and what it explains. Where a comparison can be made by shared values applied to shared evidence and interpreted claims, theory choice is rational in the sense that the community's choice can be supported by reasons available to both sides.
\end{proof}

SCI-05-a reaches a related conclusion from the side of method: the norms that unite the scientific methods, answerability, public checkability, proportion and severity, are not paradigm-relative, and so are part of what keeps theory choice across revolutions rational. SCI-06-a argues that successive paradigms are partially translatable and that theory choice remains rational across revolutions; the present argument shows that Kuhn's later account of translation already contains the premises for that conclusion.

\section{Objections and replies}\label{s:objections}

\begin{objection}[Holism]
The terms of a cluster are learned together and depend on one another. If even one term is untranslatable, the theory's laws, which contain it, are untranslatable, and the laws are what is compared.
\end{objection}

\begin{reply}
The laws are compared through their consequences, and the consequences that bear on choice are stated largely in translatable vocabulary: the weight gained by a metal on calcination, the volume of gas absorbed. That the Lavoisierian explanation of the weight gain cannot be translated into phlogistic terms does not prevent both parties from recognizing that one theory predicts the gain and the other must accommodate it.
\end{reply}

\begin{objection}[Davidson]
On Davidson's argument, the very idea of an untranslatable language is incoherent, since we could not recognize as a language what we cannot translate \citep{Davidson1974}. Kuhn's local thesis is therefore false, not merely compatible with rational choice.
\end{objection}

\begin{reply}
Davidson's argument targets total untranslatability, and it is consistent with Kuhn's local thesis: we recognize the phlogistic chemist as speaking a language because most of it is translatable, which is the background against which local failures can be identified. The argument and Kuhn's account converge on partial translatability.
\end{reply}

\section{Conclusion}\label{s:conclusion}

Kuhn's later account confines untranslatability to a local cluster, concedes that the cluster can be interpreted, and grounds theory choice in values shared across revolutions. Each concession is a premise for the conclusion of SCI-06: successive paradigms are partially translatable, and theory choice across revolutions remains rational.

\begin{thebibliography}{9}
\bibitem[Davidson(1974)]{Davidson1974} Davidson, D. (1974). On the very idea of a conceptual scheme. \emph{Proceedings and Addresses of the American Philosophical Association}, 47, 5--20.
\bibitem[Kitcher(1978)]{Kitcher1978} Kitcher, P. (1978). Theories, theorists and theoretical change. \emph{The Philosophical Review}, 87(4), 519--547.
\bibitem[Kuhn(1970)]{Kuhn1970} Kuhn, T.~S. (1970). \emph{The Structure of Scientific Revolutions}. 2nd edition. University of Chicago Press.
\bibitem[Kuhn(1977)]{Kuhn1977} Kuhn, T.~S. (1977). Objectivity, value judgment, and theory choice. In \emph{The Essential Tension: Selected Studies in Scientific Tradition and Change}, 320--339. University of Chicago Press.
\bibitem[Kuhn(1983)]{Kuhn1983} Kuhn, T.~S. (1983). Commensurability, comparability, communicability. In P.~D. Asquith and T.~Nickles (eds.), \emph{PSA 1982}, vol.~2, 669--688. Philosophy of Science Association.
\end{thebibliography}

\end{document}
